\documentclass[a4paper]{article}

\usepackage{kdiss_conference_autumn}% autumn conference proceeding style file
\usepackage{amsmath,graphicx,chicago}
\usepackage{amssymb}
\usepackage{multirow}
\usepackage{array}
\usepackage{hhline}

%% ==================== 여기는 고치지 마십시오. =======================
\submit{final}
\thefirstpage{1}
\thelastpage{1}
%% ====================================================================

\begin{document}

\title{Volatility Forecasting for Nonstationary Time Series: A Performance Comparison of Sequential and Parallel Algorithms}

\eauthor{Taejun Yoon$^{*}$~$\cdot$~Chaewon Kim$^{*}$~$\cdot$~Nakhoon Son$^{*}$}

\eaddress{${}^{*}$Department of Statistics, Keimyung University, Daegu, Korea}

\begin{eabstract}
Cryptocurrency prices form a nonstationary time series whose level does not revert and whose volatility stays high or low for long periods before shifting abruptly.
Selecting a forecasting model by a single average score can therefore choose an unsuitable model exactly when the market moves the most.
Using 15-minute prices of 20 assets traded on Upbit from October 2023 to October 2026, we forecast realized volatility from 15 minutes to 12 hours ahead and compare 27 models:
sequential models that read past observations one step at a time (GARCH and recurrent neural networks) and parallel models that process the whole input window at once
(tree ensembles, kernel regression, Transformers, and pretrained time-series foundation models).
Forecast times are divided into five regimes by the volatility observed just before the forecast, and differences in forecast loss are tested for statistical significance.
Tree ensembles and kernel regression give the smallest loss in calm and moderate regimes, whereas recurrent neural networks give the smallest loss for horizons of one hour or longer in the most volatile regime,
although their difference from tree and kernel models is not statistically significant.
In that regime every parallel deep learning model is significantly worse than the best model, and recurrent neural networks keep their advantage over parallel deep learning models even when both receive the same input.
These results suggest that volatility forecasting for nonstationary series should select models according to the volatility regime rather than search for a single best model.
\end{eabstract}

\ekeywords{Nonstationary time series, volatility forecasting, recurrent neural network, Transformer, cryptocurrency.}

\end{document}
